% --- INICIO DE PLANTILLA PERSONALIZADA ---
\documentclass[oneside,11pt]{Latex/Classes/PhDthesisPSnPDF}

% --- Paquetes Originales ---
\usepackage{blindtext}
\usepackage{actuarialangle} 
\usepackage{tabularx}
\usepackage{float}
\usepackage{multirow, array}
\usepackage[usenames,dvipsnames,svgnames,table]{xcolor}
\usepackage{longtable}
\usepackage{latexsym,amsmath,amssymb,amsfonts}
\usepackage{enumitem}
\usepackage{xparse}
\usepackage{booktabs}
\usepackage[round, sort, numbers]{natbib}  
\usepackage{adjustbox}
\usepackage[utf8]{inputenc}
\usepackage{caption}
\usepackage{graphicx}
\usepackage{hyperref}
\usepackage{cleveref}

% Definición de colores
\definecolor{miblue}{RGB}{68, 114, 196}

% --- CAMBIO CRÍTICO 1: Usar \input para macros en el preámbulo ---
% \include da problemas antes del begin document. \input es seguro.
\input{Latex/Macros/MacroFile1}

% --- Definiciones de seguridad (por si la clase falla al cargar alguna variable) ---
\providecommand{\facultad}[1]{\def\lafacultad{#1}}
\providecommand{\escudofacultad}[1]{\def\elescudo{#1}}
\providecommand{\degree}[1]{\def\elgrado{#1}}
\providecommand{\director}[1]{\def\eldirector{#1}}
\providecommand{\lugar}[1]{\def\ellugar{#1}}
\providecommand{\degreedate}[1]{\def\lafecha{#1}}

% --- Configuración para bloques de código de R (Pandoc) ---

%%%%%%%%%%%%%%%%%%%%%%%%%%%%%%%%%%%%%%%%%%%%%%%%%%%%%%%%%%%%%%%%%%%%%%%%%%%%%%%%
%                         DATOS (Conectados a RMarkdown)                       %
%%%%%%%%%%%%%%%%%%%%%%%%%%%%%%%%%%%%%%%%%%%%%%%%%%%%%%%%%%%%%%%%%%%%%%%%%%%%%%%%
\title{Analisis Estadístico de la temporada de la NBA 2019}
\author{Samuel Ruiz \\ Scarlet Casique}

% Lógica para Facultad: Si la defines en el YAML la usa, si no, usa la default

\facultad{Facultad de Ciencias Económicas y Sociales \\ Escuela de
Estadistica y Ciencias Actuariales}
  \escudofacultad{Latex/Classes/Escudos/faces}

\degree{Computación I}
\director{Jesus Ochoa \\ Oliver Triveño}
\degreedate{Abril 2026} 
\lugar{Caracas}

\portadatrue 

% Metadatos PDF
\keywords{}
\subject{}

% --- CORRECCIÓN PARA PANDOC NUEVO (Imágenes) ---
\makeatletter
\providecommand{\pandocbounded}[1]{#1}
\makeatother


%%%%%%%%%%%%%%%%%%%%%%%%%%%%%%%%%%%%%%%%%%%%%%%%%%%%%
%                   DOCUMENTO                       %
%%%%%%%%%%%%%%%%%%%%%%%%%%%%%%%%%%%%%%%%%%%%%%%%%%%%%
\begin{document}

\maketitle

%%%%%%%%%%%%%%%%%%%%%%%%%%%%%%%%%%%%%%%%%%%%%%%%%%%%%
%                  PRÓLOGO                          %
%%%%%%%%%%%%%%%%%%%%%%%%%%%%%%%%%%%%%%%%%%%%%%%%%%%%%
\frontmatter

% Puedes agregar dedicatorias desde el YAML si quieres, o dejarlas fijas aquí

%%%%%%%%%%%%%%%%%%%%%%%%%%%%%%%%%%%%%%%%%%%%%%%%%%%%%
%                   ÍNDICES                         %
%%%%%%%%%%%%%%%%%%%%%%%%%%%%%%%%%%%%%%%%%%%%%%%%%%%%%
\setcounter{secnumdepth}{3}
\setcounter{tocdepth}{3}

\tableofcontents



\cleardoublepage

  \cleardoublepage       % Empieza en página derecha (estándar en tesis)
  \chapter*{Resumen}     % Crea el título "Resumen" sin numeración (Capítulo *)
  \addcontentsline{toc}{chapter}{Resumen} % (Opcional) Lo añade al índice
  
  En el presente trabajo, se analiza el impacto del factor de localía en
la NBA durante la temporada de 2019, utilizando metodos de estadística
descriptiva (medidas de tendencia central, medidas de dispersion y
correlacion), para determinar si los equipos mejoran su rendimiento,
jugando en la localia, a diferencia de cuando juegan de
visitante.             % Aquí se pega el texto que escribiste en el YAML
  
  \cleardoublepage

%%%%%%%%%%%%%%%%%%%%%%%%%%%%%%%%%%%%%%%%%%%%%%%%%%%%%
%                   CONTENIDO (El Corazón)          %
%%%%%%%%%%%%%%%%%%%%%%%%%%%%%%%%%%%%%%%%%%%%%%%%%%%%%
\mainmatter
\def\baselinestretch{1.5}

% --- CAMBIO CRÍTICO 2: Aquí RMarkdown inyecta todo tu texto ---
\chapter*{Introducción}

En la NBA comunmente reconocen el factor de ``localía'' y dicho factor
lo asocian a una gran frecuencia de victorias en los partidos. No
obstante, independientemente del resultado final, existe la
interrogante, sobre si éste ``factor'' o ``condición'' se puede traducir
como una mejora real en el rendimiento estadístico de los equipos.

Por ende, el presente trabajo analiza la temporada 2019 con el objetivo
de cuantificar el impacto de la localía, enfocandose en métricas claves
(puntos, asistencias, rebotes, etc\ldots). El mismo se realizará,
mediante el uso de medidas de tendencia central, medidas de dispersión,
diagramas de cajas, y correlaciones. Para determinar si los equipos
rinden mejor como locales, si la eficiencia de tiro explica sus
victorias siendo locales y tambien, para observar cual estadística se
asocia fuertemente con el triunfo estando en dicha ``condición''.

\chapter{Planteamiento del Problema}

Si bien el fenómeno de la `ventaja de localía' es un factor reconocido
en el baloncesto profesional por su impacto en los resultados finales,
persiste una falta de profundidad en el análisis de las variables que lo
componen.

Actualmente, no se ha determinado con precisión si esta superioridad se
refleja de manera uniforme en las métricas de rendimiento
técnico-táctico de los equipos. Por consiguiente, el presente estudio
aborda la ausencia de una caracterización estadística detallada que
permita cuantificar el impacto real de la localía por equipo y su
correlación directa con la eficacia competitiva durante la temporada
2019 de la NBA.

Argumento final. Debido a esto, se plantea las siguientes preguntas de
investigación:\\

\textbf{1.} ¿Existe una diferencia estadísticamente significativa entre
el porcentaje de victorias como ``local'', de los equipos de la NBA y su
eficiencia de tiros de campo (FG\_PCT), tiros de tres puntos (FG3\_PCT)
y tiros libres (FT\_PCT) cuando juegan en casa de frente a cuando juegan
de visitante durante la temporada de 2019?

\textbf{2.} ¿Los equipos en promedio anotan más puntos, lanzan con mejor
porcentaje (en tiros de campo, triples y libres), dan más asistencias y
capturan más rebotes, cuando juegan de locales, a diferencia de cuando
juegan de visitantes?

\textbf{3.} En base a las estadísticas de los equipos en sus partidos
como local (los puntos, las asistencias, y los rebotes) ¿cuál tiene una
correlación más alta con la victoria?

\section{Justificación}

Este trabajo se realiza debido a la necesidad de aplicar técnicas de
estadística descriptiva junto a herramientas computacionales, para
cuantificar, de manera objetiva, el impacto real de la ventaja de la
localía en el rendimiento de los equipos de la NBA. Ademas este trabajo
ayudaría a entrenadores o aficionados, a identificar cuales aspectos
específicos del juego, favorecen más a los equipos, al estar jugando en
casa. Cabe mencionar, que este estudio tambien aporta una metodología
replicable, haciendo uso de R, para el análisis exploratorio y
descriptivo de grandes volúmenes de información.

\section{Objetivos}

\subsection{Objetivo General}
\begin{itemize}
  \item{Analizar el impacto de la localía en el rendimiento estadístico y el resultado final de los equipos de la NBA durante la temporada 2019, comparando los factores de anotación y eficiencia entre partidos de locales y visitantes.}
\end{itemize}
\subsection{Objetivos Específicos}
\begin{enumerate}
  \item Determinar, para cada equipo de la NBA, su porcentaje de victorias como local durante la temporada 2019, y analizar la distribución de las victorias dado el tipo de lanzamiento para evaluar si la ventaja de localia se refleja en el resultado final del marcador por una mayor eficiencia de lanzamiento como local.
  \item Calcular, para cada equipo, la diferencia en su rendimiento promedio como local frente a su rendimiento promedio como visitante en las siguientes estadísticas: puntos anotados, porcentaje de tiros de campo, porcentaje de tiros de tres puntos, porcentaje de tiros libres, asistencias y rebotes totales. Analizar si estas diferencias son, en promedio, favorables a la condición de local.
  \item Analizar la correlación entre las principales estadísticas de rendimiento (puntos, asistencias y rebotes) de los equipos en sus partidos como local y la victoria en dichos partidos, para identificar qué variable presenta una asociación más fuerte con el éxito en condición de local.
\end{enumerate}

\section{Cobertura}

\subsection{Horizontal}

Describir la cobertura Horizontal del trabajo de investigación\\

\begin{table}[ht]
\centering
\caption{Variables consideradas} 
\resizebox{11cm}{!} {
\begin{tabular}{ccc}
  \hline
\ \ \ \ \ NOMBRE DE LA VARIABLE \\ 
  \hline
   VARIABLE 1 \\ 
VARIABLE 2\\ 
VARIABLE n\\ 
   \hline
\end{tabular}
}
\end{table}

\subsection{Vertical}

La cobertura del trabajo de investigación, toma un total de 1.242
observaciones (o registros) en base al dataset original. Cabe destacar
que cada fila representa un partido único en la NBA, y el alcance del
dataset, abarca desde la temporada 2013 hasta la 2022, en lo que en el
mismo se puede analizar de manera múltiple; solo que en la presente
investigación tomaremos la temporada 2019

\section{Periodo de Referencia}

La fuente de los datos proviene del repositorio de Kaggle. El periodo de
referencia competente para el analisis descriptivo, se centrara en la
temporada 2019; donde las consistencia de los datos es mejor. Y el
procesamiento de la información se realizó con el corte en abril de 2026

\section{Variables}

La variable de ``interés'' en el estudio es (HOME\_TEAM\_WINS), la cual
representa el éxito competitivo por parte del equipo local

\chapter{Marco Teórico}

\section{Bases Teóricas}

Resumen de las bases teoricas explicadas en el capitulo y su
importancia\\

\subsection{Definición 1}

Descripción de la definición 1\\

\subsection{Definición 2}

Descripción de la definición 2\\

\subsection{Definición n}

Descripción de la definición n\\

\chapter{Marco Metódico}

En esta sección se presentan las metodologías y/o test necesarios
relacionados con la práctica divulgada en el marco teórico, así como las
fuentes de datos.

\section{Bases metódicas utilizadas en el trabajo de investigación}

Desarrollo\\

\subsection{Item 1}

Descripción\\

\subsection{Item 2}

Descripción\\

\subsubsection{Item 2.2}

Descripción\\

\subsection{Item n}

Descripción\\

\chapter{Análisis de Resultados}

\{\small Este capítulo presenta los resultados obtenidos al aplicar las
metodologías descritas en el capítulo anterior, así como del análisis
estadístico descriptivo de las variables estudiadas y las ventajas y
desventajas al aplicar dichos tests.\}

\section{Presentación de resultados 1}

Descripción\\

\section{Presentación de resultados 2}

Descripción\\

\section{Presentación de resultados n}

Descripción\\

\chapter{Conclusiones y Recomendaciones}

\textbf{1.-} Conclusion 1 del trabajo de investigación.\\

\textbf{Recomendación} Recomendación 1 del trabajo de investigación.\\

\textbf{2.-} Conclusion 2 del trabajo de investigación.\\

\textbf{Recomendación} Recomendación 2 del trabajo de investigación.\\

\textbf{n.-} Conclusion n del trabajo de investigación.\\

\textbf{Recomendación} Recomendación n del trabajo de investigación.\\

\appendix
\renewcommand{\thechapter}{\Alph{chapter}}

\chapter{Anexo A}

\chapter{Anexo B}

\chapter{Anexo C}

%%%%%%%%%%%%%%%%%%%%%%%%%%%%%%%%%%%%%%%%%%%%%%%%%%%%%
%                   REFERENCIAS                     %
%%%%%%%%%%%%%%%%%%%%%%%%%%%%%%%%%%%%%%%%%%%%%%%%%%%%%
% Mantenemos la estructura natbib de tu plantilla
\renewcommand{\bibname}{Lista de Referencias}
\bibliographystyle{apalike} 
\bibliography{bibliografia.bib} 

%%%%%%%%%%%%%%%%%%%%%%%%%%%%%%%%%%%%%%%%%%%%%%%%%%%%%
%                   APÉNDICES                       %
%%%%%%%%%%%%%%%%%%%%%%%%%%%%%%%%%%%%%%%%%%%%%%%%%%%%%

\end{document}
